% A visual introduction with full core grammars, followed by reference slides.
% Build: latexmk -pdf -interaction=nonstopmode -halt-on-error dot-mnf-fcdot.tex
% Presenter commentary and Lean references are in speaker-notes.md.
\documentclass[aspectratio=169,12pt]{beamer}
\usepackage[T1]{fontenc}
\usepackage{lmodern,amsmath,amssymb,mathtools,stmaryrd,booktabs}
\usefonttheme{professionalfonts}
\definecolor{ink}{HTML}{182B3A}
\definecolor{source}{HTML}{225FA0}
\definecolor{target}{HTML}{A34428}
\definecolor{muted}{HTML}{53616B}
\setbeamercolor{normal text}{fg=ink,bg=white}
\setbeamercolor{frametitle}{fg=ink}
\setbeamercolor{structure}{fg=ink}
\setbeamerfont{frametitle}{size=\large,series=\bfseries}
\setbeamerfont{footline}{size=\scriptsize}
\setbeamersize{text margin left=10mm,text margin right=10mm}
\setbeamertemplate{navigation symbols}{}
\setbeamertemplate{headline}{}
\setbeamertemplate{frametitle}{\vspace{4mm}\insertframetitle\par\vspace{2mm}}
\setbeamertemplate{footline}{\hfill\insertframenumber\hspace{10mm}\vspace{3mm}}
\title{DOT-MNF to FCdot}
\author{}
\date{}
\pdfstringdefDisableCommands{\def\translate#1{#1}}
\hypersetup{pdftitle={DOT-MNF to FCdot},pdfauthor={},
  pdfsubject={The calculi, translation, and safety for WadlerFest DOT},
  colorlinks=true,urlcolor=source}

\newcommand{\src}[1]{{\color{source}#1}}
\newcommand{\tgt}[1]{{\color{target}#1}}
% The brackets belong to the target; their argument belongs to the source.
\newcommand{\tr}[1]{\tgt{\llbracket\src{#1}\rrbracket}}
\newcommand{\tel}[1]{\tgt{\mathsf{tel}(\src{#1})}}
\newcommand{\tels}[2]{\tgt{\mathsf{tel}_{\src{#1}}(\src{#2})}}
\newcommand{\ann}[1]{\src{\mathsf{strip}(#1)}}
\newcommand{\er}[1]{{\color{ink}\lfloor#1\rfloor}}
\newcommand{\sub}{\mathrel{<:}}
\newcommand{\mem}[3]{\{#1:#2\mathinner{.\!.}#3\}}
\newcommand{\fld}[2]{\{#1:#2\}}
\newcommand{\obj}[2]{\mu(#1.\,#2)}
\newcommand{\nm}[2]{#1\mathbin{\cdot}#2}
\newcommand{\cast}[2]{\mathsf{cast}(#1,#2)}
\newcommand{\proj}[3]{\mathsf{proj}(#1,#2,#3)}
\newcommand{\has}[1]{\mathsf{has}\;#1}
\newcommand{\bnd}[1]{\mathsf{bnd}\;#1}
\newcommand{\rt}[1]{\mathsf{root}(#1)}
\newcommand{\readp}[2]{\mathsf{read}(#1,#2)}
\newcommand{\kw}[1]{\mathsf{#1}}
\newcommand{\gap}{\par\vspace{3mm}}
\newcommand{\captionline}[1]{\par\vspace{2mm}{\small\color{muted}#1}\par}

\begin{document}

\begin{frame}{DOT-MNF to FCdot}
  \vspace{2mm}
  {\Large Type safety through explicit evidence}
  \vspace{7mm}
  \[
  \begin{array}{c@{\quad}c@{\quad}c}
    \src{\mathcal{H}\mathrel{::}\Gamma\vdash t:T}&\longmapsto&\tgt{\tr{\Gamma}\vdash\tr{\mathcal{H}}:\tr{T}}\\[2mm]
    \text{\src{implicit subtyping}}&&\text{\tgt{explicit casts}}\\[5mm]
    \multicolumn{3}{c}{\er{\tr{\mathcal{H}}}=\er{\src{t}}}
  \end{array}
  \]
  \gap
  Translate a typing derivation \src{$\mathcal{H}$}. Reuse the target's safety proof.
  \captionline{Erasure $\er{\,\cdot\,}$ removes types and evidence.\\
  The result covers the annotated WadlerFest rules and every retained-let reduction order.\\
  \src{Blue}: source.\quad\tgt{Rust}: target.}
\end{frame}

\begin{frame}{DOT: fields and abstract type members}
  \begin{center}
  \begin{tabular}{@{}ll@{}}
    $\src{\mem{A}{S}{T}}$ & type member with bounds $\src{S}$ and $\src{T}$\\[2mm]
    $\src{\fld{a}{T}}$ & field of type $\src{T}$\\[2mm]
    $\src{x.A}$ & the type member of variable $\src{x}$
  \end{tabular}
  \end{center}
  \gap
  \textbf{Example}, assuming $\src{n:R}$:
  \[
    \src{\kw{let}\ x=\underbrace{\nu(z.a=n)}_{\text{object}}\ \kw{in}\
      \underbrace{x.a}_{\text{projection}}:R}
  \]
  \captionline{$\src{z}$ names self.
  MNF: applications and projections use variables; a \texttt{let} names a computation.}
\end{frame}

\begin{frame}{DOT-MNF syntax: types}
  \textbf{Types}
  \[
  \src{\begin{aligned}
    S,T&::=\top\mid\bot\mid\mem{A}{S}{T}\mid\fld{a}{T}\mid x.A\\[2mm]
       &\quad\mid\forall(x:S)T\mid\mu(z.T)\mid S\wedge T
  \end{aligned}}
  \]
  \gap
  \textbf{Recursive introduction and elimination}
  \[
    \src{\frac{\Gamma\vdash x:T[x/z]}{\Gamma\vdash x:\mu(z.T)}}
    \qquad
    \src{\frac{\Gamma\vdash x:\mu(z.T)}{\Gamma\vdash x:T[x/z]}}
  \]
  \captionline{The body $\src{T}$ is arbitrary: selections, functions,
  intersections, and nested recursion are all admitted.
  $\src{A,B}$ are type labels; $\src{a,b}$ are field labels.}
\end{frame}

\begin{frame}{DOT-MNF syntax: terms and definitions}
  \[
  \begin{array}{r@{\;\src{::=}\;}l@{\qquad}l}
    \src{t,u}&\src{x\mid v\mid x\,y\mid x.a\mid\kw{let}\ x=t\ \kw{in}\ u}
      &\text{terms}\\[5mm]
    \src{v}&\src{\lambda(x:S).t\mid\nu(z.d)}&\text{values}\\[5mm]
    \src{d}&\src{A=T\mid a=t\mid d\wedge d}&\text{definitions}
  \end{array}
  \]
  \gap
  Applications and projections use variables. Field bodies may compute.
  \captionline{The object binds self $\src{z}$ in all its definitions.
  Well-typed literals have distinct labels.}
\end{frame}

\begin{frame}{The annotated WadlerFest frontend}
  The published object syntax records its self type:
  \[
    \src{v::=\lambda(x:S).t\mid\nu(z:T)d}.
  \]
  Types, terms, and definitions otherwise have the preceding grammars.
  \gap
  Its object rule checks definitions under an \textbf{opened self assumption}:
  \[
    \src{\frac{\Gamma,z:T\vdash d:T}
      {\Gamma\vdash\nu(z:T)d:\mu(z.T)}}.
  \]
  \captionline{Here $\src{T}$ may mention $\src{z}$.
  Definition typing requires disjoint labels across conjunctions.
  The public frontend enforces distinct type-label and field-label categories.}
\end{frame}

\begin{frame}{Opened self is represented by recursive elimination}
  DOT-MNF removes the annotation and remembers the literal at self:
  \[
    \ann{\nu(z:T)d}=\src{\nu(z.\kw{strip}(d))}.
  \]
  The local self assumption is recursive. One rule recovers the opened type:
  \[
    \src{z:\mu(z.T):=d\quad\Longrightarrow\quad z:T}.
  \]
  \gap
  \textbf{Typing correspondence.} Every closed annotated derivation yields
  \[
    \src{\varnothing\vdash_{\!W}t:T}
      \quad\Longrightarrow\quad
    \src{\varnothing\vdash_{\!D}\kw{strip}(t):T}.
  \]
  \captionline{$\src{W}$ denotes the annotated rules, $\src{D}$ DOT-MNF.
  No good-bounds or inertness premise is needed.}
\end{frame}

\begin{frame}{Bad bounds hide a term's runtime shape}
  \[
    \src{\frac{\Gamma\vdash x:\mem{A}{S}{T}}
      {\Gamma\vdash S\sub x.A\qquad\Gamma\vdash x.A\sub T}}
  \]
  \gap
  For a hypothetical parameter $\src{x:\mem{A}{\top}{\bot}}$:
  \[
    \src{\underbrace{U\sub\top}_{\text{any }U}
      \sub x.A\sub
      \underbrace{\bot\sub V}_{\text{any }V}}
  \]
  \gap
  Subsumption can now assign any type to an already typed term.
  \captionline{These assumptions are legal under a lambda.
  A typed runtime store cannot realize them.}
\end{frame}

\begin{frame}{FCdot syntax: types and propositions}
  \textbf{Types and telescopes}
  \[
  \tgt{\begin{aligned}
    \tau,\sigma&::=\bot\mid\nm{x}{\ell}\mid\Pi(x:\tau)\sigma\mid\obj{z}{\Phi}\\[2mm]
    \Phi,\Psi&::=[\,]\mid\Phi,P
  \end{aligned}}
  \]
  \gap
  \textbf{Propositions about self}
  \[
    \tgt{P::=\tau\sqsubseteq\sigma\mid\tau\doteq\sigma
      \mid\has{a}\mid\bnd{\tau}}
  \]
  \captionline{$\tgt{x,y,z}$ are variables.
  $\tgt{\obj{z}{\Phi}}$ binds $\tgt{z}$ throughout $\tgt{\Phi}$.
  $\tgt{\top:=\obj{z}{[\,]}}$.\\
  $\tgt{\bnd{\tau}}$ includes self in $\tgt{\tau}$.
  Its type may mention self; recursive unfolding will open that dependency.}
\end{frame}

\begin{frame}{FCdot: member type names}
  $\tgt{\nm{x}{\ell}}$ is a type name attached to variable $\tgt{x}$ and label $\tgt{\ell}$.
  \gap
  \[
  \begin{array}{l@{\qquad}l}
    \tgt{\nm{x}{A}}&\text{the type named by member }\tgt{A}\\[4mm]
    \tgt{\nm{x}{a}}&\text{the result type of field }\tgt{a}
  \end{array}
  \]
  \gap
  \[
    \tr{x.A}=\tgt{\nm{x}{A}}.
  \]
  \captionline{An object gives these names bounds or exact definitions.
  Field presence is a separate fact.}
\end{frame}

\begin{frame}{FCdot: member names and object facts}
  The self variable $\tgt{z}$ can occur in the types inside $\tgt{P}$:
  \gap
  \[
    \tgt{\nm{z}{a}}\quad\text{is an instance of the type grammar's }\tgt{\nm{x}{\ell}}.
  \]
  \[
    \tgt{\obj{z}{[\underbrace{\has{a}}_{\text{self has field }a},\
      \underbrace{\nm{z}{a}\sqsubseteq\tau}_{\text{bound on its type}}]}}
  \]
  \captionline{Inclusion and equality mention self through their types.
  $\tgt{\has{a}}$ and $\tgt{\bnd{\tau}}$ refer to the enclosing self implicitly.}
\end{frame}

\begin{frame}{Evidence, judgments, and derivations}
  \begin{center}
  \begin{tabular}{@{}ll@{}}
    Proposition syntax & $\tgt{\tau\sqsubseteq\tau}$\\[5mm]
    Proof-term definition & $\tgt{e:=\kw{refl}(\tau)}$\\[5mm]
    Typing judgment & $\tgt{\Delta\vdash e:\tau\leq\tau}$\\[5mm]
    Derivation of that judgment
      & $\tgt{\mathcal{E}\mathrel{::}\Delta\vdash e:\tau\leq\tau}$
  \end{tabular}
  \end{center}
  \gap
  A cast contains the proof term $\tgt{e}$.
  \captionline{The derivation $\tgt{\mathcal{E}}$ establishes its typing judgment.
  A telescope contains proposition syntax.}
\end{frame}

\begin{frame}{Each use of subsumption becomes an explicit cast}
  \[
    \src{\frac{\Gamma\vdash t:S\qquad\Gamma\vdash S\sub T}
      {\Gamma\vdash t:T}}
  \]
  \vspace{3mm}
  \begin{center}The subtyping derivation translates to a proof term $\tgt{e}$\end{center}
  \vspace{3mm}
  \[
    \tgt{\frac{\Delta\vdash u:\tau\qquad
        \Delta\vdash e:\tau\leq\sigma}
      {\Delta\vdash\cast{u}{e}:\sigma}}
  \]
  \captionline{$\tgt{\Delta}$ is the target context.
  The premise $\tgt{\Delta\vdash e:\tau\leq\sigma}$ checks its evidence argument.}
\end{frame}

\begin{frame}{FCdot syntax: atoms}
  An atom is a variable under wrappers:
  \[
  \tgt{\begin{aligned}
    r,s&::=x\mid\cast{r}{e}\\[2mm]
       &\quad\mid\kw{fold}(z.\Phi,r)\mid\kw{unfold}(r)\\[2mm]
       &\quad\mid\kw{both}(z.\Phi,z.\Psi,r,s)
  \end{aligned}}
  \]
  \captionline{Fold/unfold introduce/eliminate a recursive self type.
  Both combines two views of the same variable.}
\end{frame}

\begin{frame}{The root of an atom: all cases}
  \[
  \tgt{\begin{aligned}
    \rt{x}&=x\\[3mm]
    \rt{\cast{r}{e}}&=\rt{r}\\[3mm]
    \rt{\kw{fold}(z.\Phi,r)}&=\rt{r}\\[3mm]
    \rt{\kw{unfold}(r)}&=\rt{r}\\[3mm]
    \rt{\kw{both}(z.\Phi,z.\Psi,r,s)}&=\rt{r}
  \end{aligned}}
  \]
  \gap
  \captionline{Root traverses atom wrappers; it does not inspect object definitions.\\
  For $\tgt{\kw{both}}$, it uses the first operand; typing requires $\tgt{\rt{r}=\rt{s}}$.}
\end{frame}

\begin{frame}{FCdot syntax: terms and values}
  \textbf{Terms}
  \[
  \tgt{\begin{aligned}
    u,w&::=r\mid v\mid r\,s\mid\proj{r}{a}{p}\\[2mm]
       &\quad\mid\kw{let}\ x=u\ \kw{in}\ w\mid\cast{u}{e}
  \end{aligned}}
  \]
  \gap
  \textbf{Values}
  \[
    \tgt{v::=\lambda(x:\tau).u\mid\kw{obj}(z;W;F)\mid\cast{v}{e}}
  \]
  \captionline{Applications and projections use atoms.
  $\tgt{e}$ is inclusion evidence; $\tgt{p}$ is field-presence evidence.
  A cast value remains a value.}
\end{frame}

\begin{frame}{FCdot syntax: object definitions}
  \[
  \begin{array}{r@{\;\tgt{::=}\;}l@{\qquad}l}
    \tgt{W}&\tgt{[\,]\mid W,\ell=\tau}&\text{type witnesses}\\[4mm]
    \tgt{F}&\tgt{[\,]\mid F,a=u}&\text{field bodies}
  \end{array}
  \]
  \gap
  The object $\tgt{\kw{obj}(z;W;F)}$ binds self in both lists.
  \gap
  \[
    \tgt{W=[a=\tau]},\qquad\tgt{F=[a=u]}.
  \]
  \captionline{Here, $\tgt{\nm{z}{a}}$ names $\tgt{\tau}$ and the field body is $\tgt{u}$.}
\end{frame}

\begin{frame}{Declarations translate to the facts they provide}
  \[
  \begin{aligned}
    \tr{\mem{A}{S}{T}}
      &=\tgt{\obj{z}{[
         \underbrace{\tr{S}\sqsubseteq\nm{z}{A}}_{\text{lower bound}},\
         \underbrace{\nm{z}{A}\sqsubseteq\tr{T}}_{\text{upper bound}}]}}
         \\[9mm]
    \tr{\fld{a}{T}}
      &=\tgt{\obj{z}{[
         \underbrace{\has{a}}_{\text{presence}},\
         \underbrace{\nm{z}{a}\sqsubseteq\tr{T}}_{\text{result type}}]}}
  \end{aligned}
  \]
  \vspace{6mm}
  Reading a fact replaces self $\tgt{z}$ by the receiver's variable.
\end{frame}

\begin{frame}{The remaining type constructors}
  \[
  \begin{aligned}
    \tr{\top}&=\tgt{\obj{z}{[\,]}},\qquad
      \tr{\bot}=\tgt{\bot},\qquad \tr{x.A}=\tgt{\nm{x}{A}}\\[5mm]
    \tr{\forall(x:S)T}&=\tgt{\Pi(x:\tr{S})\tr{T}}\\[5mm]
    \tr{S\wedge T}&=\tgt{\obj{z}{\tel{S}\mathbin{+\!+}\tel{T}}}\\[5mm]
    \tr{\mu(z.T)}&=\tgt{\obj{z}{\tels{z}{T}}}
  \end{aligned}
  \]
  \captionline{$\tel{T}$ describes a type under a fresh self binder.
  $\tels{z}{T}$ describes a body whose self is already $\src{z}$.}
\end{frame}

\begin{frame}{Intersections include arbitrary operands}
  Declarations contribute their member facts. An operand with another shape
  contributes a bound on the object itself:
  \[
    \tel{x.A}=\tgt{[\bnd{\nm{x}{A}}]}.
  \]
  \[
    \tr{x.A\wedge\fld{a}{T}}
      =\tgt{\obj{z}{[\bnd{\nm{x}{A}},\has{a},\nm{z}{a}\sqsubseteq\tr{T}]}}.
  \]
  \gap
  The first fact supports projection back to $\tgt{\nm{x}{A}}$.
  \captionline{Functions, bottom, and recursive types with nondeclaration bodies
  use the same one-bound representation as operands.\
  Every bound in $\tel{T}$ is independent of its fresh self.}
\end{frame}

\begin{frame}{Arbitrary recursive bodies retain their self dependency}
  A recursive body can be a bare self selection:
  \[
    \tr{\mu(z.z.A)}=\tgt{\obj{z}{[\bnd{\nm{z}{A}}]}}.
  \]
  For an atom $\tgt{r}$ of this type with $\tgt{\rt{r}=x}$:
  \[
  \tgt{\begin{array}{c}
    r:\obj{z}{[\bnd{\nm{z}{A}}]}\\[1mm]
    \downarrow\ \kw{unfold}\\[1mm]
    \kw{unfold}(r):\obj{w}{[\bnd{\nm{x}{A}}]}\\[1mm]
    \downarrow\ \text{cast through the bound}\\[1mm]
    \text{result type }\nm{x}{A}
  \end{array}}
  \]
  \captionline{Opening replaces self by the root before extracting the bound.
  Recursive introduction reverses the construction.
  Both adapters preserve the root and erase to the original atom.}
\end{frame}

\begin{frame}{Replace self by the receiver's variable}
  Suppose the receiver has this type and root:
  \[
    \tgt{\Delta\vdash r:\obj{z}{[\has{a},\nm{z}{a}\sqsubseteq\tau]}},\qquad
    \tgt{\rt{r}=x}.
  \]
  \gap
  The type makes two promises about $\tgt{x}$:
  \begin{center}
  \begin{tabular}{@{}ll@{}}
    $\tgt{\has{a}}$ & field $\tgt{a}$ is available at $\tgt{x}$\\[4mm]
    $\tgt{\nm{z}{a}\sqsubseteq\tau}$
      & its named type $\tgt{\nm{x}{a}}$ is included in $\tgt{\tau[x/z]}$
  \end{tabular}
  \end{center}
  \captionline{Substitute $\tgt{x}$ for $\tgt{z}$ throughout, including inside $\tgt{\tau}$.
  No field body is evaluated here.}
\end{frame}

\begin{frame}{Reading a fact produces evidence}
  Assume $\tgt{\Delta\vdash r:\obj{z}{\Phi}}$ and $\tgt{\rt{r}=x}$, with
  \[
    \tgt{\Phi=[\has{a},\nm{z}{a}\sqsubseteq\tau]},
  \]
  \begin{center}
  \begin{tabular}{@{}l@{\hspace{9mm}}l@{}}
    \textbf{Proof-term definitions}&\textbf{Typing judgments}\\[3mm]
    $\tgt{p:=\readp{r}{0}}$&$\tgt{\Delta\vdash p:x\ni a}$\\[4mm]
    $\tgt{e:=\readp{r}{1}}$&$\tgt{\Delta\vdash e:\nm{x}{a}\leq\tau[x/z]}$
  \end{tabular}
  \end{center}
  \gap
  \[
    \tgt{\readp{r}{i}:=\kw{member}(r,\kw{refl}(\obj{z}{\Phi}),i)}.
  \]
  \captionline{Here $\tgt{\kw{refl}}$ is the identity coercion; indices start at zero.}
\end{frame}

\begin{frame}{One projection, two obligations}
  Source: $\src{\Gamma\vdash x:\fld{a}{T}}$ gives $\src{\Gamma\vdash x.a:T}$.
  \captionline{Let $\tgt{r}$ be the translated receiver, with root $\tgt{x}$.}
  \[
  \begin{array}{c@{\qquad}c}
    \tgt{\Delta\vdash p:x\ni a}&\tgt{\Delta\vdash e:\nm{x}{a}\leq\tr{T}}\\[1mm]
    \text{field presence}&\text{result conversion}
  \end{array}
  \]
  \gap
  \[
    \tgt{\Delta\vdash\cast{\underbrace{\proj{r}{a}{p}}_{\text{type }\nm{x}{a}}}{e}
      :\tr{T}}
  \]
  \captionline{Access the field at its named type, then cast to the declared type.
  These are the two facts supplied by the translated declaration.}
\end{frame}

\begin{frame}{A stronger view still refers to the same member}
  \[
    \src{S=\mem{A}{\bot}{\top}},\quad
    \src{T=\mem{A}{R}{\top}},\quad
    \src{\Gamma=x:\mem{B}{S}{T},\ w:S}.
  \]
  \begin{columns}[T,onlytextwidth]
  \begin{column}{0.45\textwidth}
    \textbf{Source}
    \[
      \src{S\sub x.B\sub T}
    \]
    \[
      \src{w:T\quad\Longrightarrow\quad R\sub w.A}.
    \]
  \end{column}
  \begin{column}{0.51\textwidth}
    \textbf{Target proof terms}
    \[
    \tgt{\begin{aligned}
      e&:=\kw{trans}(\readp{x}{0},\readp{x}{1})\\[2mm]
      r&:=\cast{w}{e}
    \end{aligned}}
    \]
    \[
      \tgt{\tr{\Gamma}\vdash\readp{r}{0}:\tr{R}\leq\nm{w}{A}}.
    \]
  \end{column}
  \end{columns}
  \gap
  Reading at the cast atom uses its root $\tgt{w}$.
  \captionline{A total translation must handle these hypothetical assumptions.
  Opening a coerced package at a fresh type name would lose this connection.}
\end{frame}

\begin{frame}{Object coercions use fixed recipes}
  Given $\tgt{\Delta\vdash e:\tau\leq\sigma}$, transform the object's facts:
  \[
  \tgt{
  \begin{array}{c@{\quad\longmapsto\quad}c}
    [\has{a},\ \nm{z}{a}\sqsubseteq\tau]
      &[\has{a},\ \nm{z}{a}\sqsubseteq\sigma]\\[4mm]
    [p_a,\ e_a]&[p_a,\ e_a;e]
  \end{array}}
  \]
  \captionline{Copy the presence proof. Compose the field-type proof with $\tgt{e}$.}
  \vspace{5mm}
  These \textbf{template morphisms} use numbered facts and evidence independent of self.
  \gap
  Their composition substitutes evidence for numbered facts.
  \captionline{$\tgt{e_a;e}$ means composition.}
\end{frame}

\begin{frame}{The store keeps exact definitions}
  For a literal with witness $\tgt{W=[a=\tau]}$ and field $\tgt{a}$:
  \[
    \tgt{\Delta\vdash\chi:
      \underbrace{\obj{z}{[\nm{z}{a}\doteq\tau,\has{a}]}}_{\text{precise literal type}}}
    \quad\leq\quad
    \tgt{\underbrace{\obj{z}{[\has{a},\nm{z}{a}\sqsubseteq\tau]}}_{\text{advertised field type}}}.
  \]
  \captionline{Equality supplies the advertised bound. Field presence stays unchanged.}
  \gap
  A \textbf{typed store} keeps bare lambdas and object literals at their own types.
  \gap
  \textbf{Inertness remains an invariant: allocation stores the literal and moves casts to its uses.}
\end{frame}

\begin{frame}{Worked example: an object containing identity}
  \[
    \src{R:=\forall(y:\top)\top},\qquad
    \tgt{\rho:=\tr{R}=\Pi(y:\top)\top}.
  \]
  \gap
  \[
  \src{\begin{aligned}
    t={}&\kw{let}\ x=\nu(z.\,a=\lambda(y:\top).y)\\[2mm]
       &\kw{in}\ x.a
  \end{aligned}}
  \]
  \[
    \src{\mathcal{H}\mathrel{::}\varnothing\vdash t:R}.
  \]
  \captionline{The program returns the identity function stored in field $\src{a}$.
  The projection opens $\src{x}$'s recursive object type.}
\end{frame}

\begin{frame}{The field body is cast to its named type}
  $\tgt{\Delta_z}$ records literal self $\tgt{z}$ with witness $\tgt{a=\rho}$.
  \begin{center}
  \begin{tabular}{@{}l@{\qquad}l@{}}
    \textbf{Proof-term definitions}&\textbf{Typing judgments}\\[3mm]
    $\tgt{\delta_z:=\kw{def}(z,a)}$
      &$\tgt{\Delta_z\vdash\delta_z:\nm{z}{a}\equiv\rho}$\\[4mm]
    $\tgt{\kappa_z:=\kw{eqToLe}(\kw{symm}(\delta_z))}$
      &$\tgt{\Delta_z\vdash\kappa_z:\rho\leq\nm{z}{a}}$
  \end{tabular}
  \end{center}
  \gap
  \[
  \tgt{\begin{aligned}
    W&=[a=\rho]\\[3mm]
    F&=[a=\cast{\lambda(y:\top).y}{\kappa_z}]
  \end{aligned}}
  \]
  \captionline{The identity has type $\tgt{\rho}$.
  Its cast gives the stored field body type $\tgt{\nm{z}{a}}$.}
\end{frame}

\begin{frame}{The literal's exact facts imply its declared type}
  \[
  \tgt{\begin{aligned}
    \Psi&=[\underbrace{\nm{z}{a}\doteq\rho}_{\text{fact 0}},\
             \underbrace{\has{a}}_{\text{fact 1}}]\\[3mm]
    \Phi&=[\underbrace{\has{a}}_{\text{fact 0}},\
             \underbrace{\nm{z}{a}\sqsubseteq\rho}_{\text{fact 1}}]
  \end{aligned}}
  \]
  \captionline{Copy presence from fact 1. Use equality at fact 0 as an upper bound.}
  \gap
  \begin{center}
  \begin{tabular}{@{}l@{\qquad}l@{}}
    \textbf{Proof terms}
      &$\tgt{m:=\kw{le}(\kw{has}(\kw{nil},1),\kw{none},\kw{eq}(0),\kw{none})}$\\[3mm]
      &$\tgt{\chi:=\kw{obj}(\Psi,m)}$\\[5mm]
    \textbf{Typing judgment}
      &$\tgt{\varnothing\vdash\chi:\obj{z}{\Psi}\leq\obj{z}{\Phi}}$
  \end{tabular}
  \end{center}
\end{frame}

\begin{frame}{The complete translated program}
  Inside the let body, abbreviate the receiver and its two proofs:
  \[
  \tgt{\begin{aligned}
    r_x&:=\kw{fold}(z.\Phi,\kw{unfold}(x))\\[1mm]
    p_x&:=\readp{r_x}{0},\qquad e_x:=\readp{r_x}{1}
  \end{aligned}}
  \]
  \gap
  \[
  \begin{aligned}
    \tgt{u}:=\tr{\mathcal{H}}=\tgt{\kw{let}\ x={}}&\tgt{\kw{cast}(\kw{obj}(z;[a=\rho];}\\[1mm]
       &\tgt{\qquad[a=\cast{\lambda(y:\top).y}{\kappa_z}]),\chi)}\\[2mm]
    \tgt{\kw{in}\quad}&\tgt{\cast{\proj{r_x}{a}{p_x}}{e_x}}
  \end{aligned}
  \]
  \[
    \tgt{\varnothing\vdash u:\rho}.
  \]
\end{frame}

\begin{frame}{The source store machine}
  A state $\src{\langle\sigma,K,t\rangle}$ contains stored values,
  pending let frames, and the current term.
  \[
  \src{\begin{aligned}
    \langle\sigma,K,\kw{let}\ x=t\ \kw{in}\ u\rangle
      &\longrightarrow\langle\sigma,K\mathbin{\triangleright}(x.u),t\rangle\\[3mm]
    \langle\sigma,K\mathbin{\triangleright}(x.u),v\rangle
      &\longrightarrow\langle\sigma[x\mapsto v],K,u\rangle\\[3mm]
    \langle\sigma,K\mathbin{\triangleright}(x.u),y\rangle
      &\longrightarrow\langle\sigma,K,u[y/x]\rangle
  \end{aligned}}
  \]
  \gap
  Application looks up a lambda. Projection looks up a field body and
  substitutes the receiver for self.
  \captionline{Final states have no pending frames and return a value or variable.
  The annotated machine retains self annotations and otherwise uses the same rules.
  Fresh-variable renaming is implicit above.}
\end{frame}

\begin{frame}{Two proofs carry different runtime invariants}
  \textbf{Direct:} retain source typing under an inert runtime context.
  \[
    \src{I_D(s_D):=\exists\Gamma,R.\
      \kw{Inert}(\Gamma)\land\kw{Typed}_D(\Gamma,s_D,R)}.
  \]
  \gap
  \textbf{Via translation:} retain a typed target state with the same erasure.
  \[
    \kw{Sim}(\src{s_D}):=\exists\tgt{s_F,\tau}.\
      \tgt{\kw{Typed}_F(s_F,\tau)}\land\er{\tgt{s_F}}=\er{\src{s_D}}.
  \]
  \gap
  Both must recover a lambda for application and a field for projection.
  \captionline{Direct baseline: \href{https://arxiv.org/pdf/1706.03814}{Rapoport et al. (2017)}.
  Its invariant is adapted schematically to store notation; only the
  translation route is mechanized here.}
\end{frame}

\begin{frame}{Direct proof: inert contexts and precise typing}
  An inert context assigns only these types:
  \[
    \src{I::=\forall(x:S)T\mid\mu(z.D_=)}
  \]
  \[
    \src{D_=::=\mem{A}{T}{T}\mid\fld{a}{T}\mid D_=\wedge D_=}.
  \]
  \captionline{Type labels are distinct. Field types and function domains may be arbitrary.}
  \gap
  \textbf{Precise variable typing} $\src{\Gamma\vdash_!x:T}$ starts at
  $\src{\Gamma(x)}$ and uses only:
  \[
    \text{lookup}\quad\longrightarrow\quad
    \text{recursive unfolding / intersection projection}.
  \]
  \captionline{It exposes the declarations actually recorded in the context.}
\end{frame}

\begin{frame}{Direct proof: recover bounds through an exact member}
  Tight selection requires a precise member with equal bounds:
  \[
    \src{\frac{\Gamma\vdash_!x:\mem{A}{R}{R}}
      {\Gamma\vdash_\#R\sub x.A\qquad\Gamma\vdash_\#x.A\sub R}}.
  \]
  \gap
  Under inert $\src{\Gamma}$, inversion of tight typing gives
  \[
    \src{\Gamma\vdash_\#x:\mem{A}{S}{T}}
    \quad\Longrightarrow\quad
    \exists\src{R}.\
    \begin{cases}
      \src{\Gamma\vdash_!x:\mem{A}{R}{R}}\\[2mm]
      \src{\Gamma\vdash_\#S\sub R,\quad\Gamma\vdash_\#R\sub T}.
    \end{cases}
  \]
  \captionline{Compose these relations to recover both general selection rules.
  This is the key selection-replacement argument.}
  \gap
  Tight typing retains transitivity and general premises under new binders.
\end{frame}

\begin{frame}{Direct proof: how the inversion theorems are used}
  For a variable in inert $\src{\Gamma}$:
  \[
  \begin{array}{c@{\quad\Longrightarrow\quad}c@{\quad\Longrightarrow\quad}c}
    \src{\Gamma\vdash x:T}&\src{\Gamma\vdash_\#x:T}&\src{\Gamma\vdash_{\#\#}x:T}\\[2mm]
    \text{general}&\text{tight}&\text{invertible}
  \end{array}
  \]
  \gap
  Invertible typing starts from precise typing, then adds introductions.
  \captionline{Eliminations and introductions are separated, so induction recovers the base.}
  \gap
  For example, field inversion yields some $\src{S}$ with
  \[
    \src{\Gamma\vdash x:\fld{a}{T}}
    \quad\Longrightarrow\quad
    \src{\Gamma\vdash_!x:\fld{a}{S}\quad\land\quad\Gamma\vdash S\sub T}.
  \]
  \captionline{Value inversion then connects the context entry to the stored value.}
\end{frame}

\begin{frame}{Head forms expose the action of a coercion}
  \[
  \tgt{\begin{aligned}
    F,G::={}&\kw{bot}\mid\kw{top}\mid\kw{id}\mid\kw{pi}(e,x.f)\\[2mm]
      &\mid\kw{obj}(\vec E)\mid\kw{bnd}(i,F)\mid\kw{into}(\vec Q)
  \end{aligned}}
  \]
  \begin{center}
  \begin{tabular}{@{}ll@{}}
    $\tgt{\kw{id}}$ & endpoints have the same resolved shape\\[2mm]
    $\tgt{\kw{pi}(e,x.f)}$ & argument and result coercions\\[2mm]
    $\tgt{\kw{obj}(\vec E)}$ & one recipe per target proposition\\[2mm]
    $\tgt{\kw{bnd}(i,F)}$ & follow bound $\tgt{i}$, then $\tgt{F}$\\[2mm]
    $\tgt{\kw{into}(\vec Q)}$ & establish an object type by bounds or routes
  \end{tabular}
  \end{center}
  \captionline{Resolution follows transparent type definitions.
  Equality-to-inclusion evidence normalizes to $\tgt{\kw{id}}$.}
\end{frame}

\begin{frame}{Local templates and object entries}
  A local template refers to a numbered source proposition:
  \[
  \tgt{\begin{aligned}
    L::={}&\kw{le}(F,H,G)\mid\kw{eq}(i,b)\mid\kw{has}(i)\\[2mm]
      &\mid\kw{copyBound}(i).
  \end{aligned}}
  \]
  Object entries add a coercion from the whole source object:
  \[
  \tgt{\begin{aligned}
    E::={}&\kw{le}(F,H,G)\mid\kw{eq}(i,b)\mid\kw{has}(i)\\[2mm]
      &\mid\kw{copyBound}(i)\mid\kw{bnd}(F)\\[2mm]
    \vec E::={}&[\,]\mid\vec E,E.
  \end{aligned}}
  \]
  \captionline{$\tgt{H}$ reads an inclusion or either direction of an equality.
  $\tgt{b}$ chooses an equality's direction.
  Copied bounds remain indexed, so a later bound elimination finds them.}
\end{frame}

\begin{frame}{Routes end at a local template}
  Pairing may obtain facts through a different object type.
  Its entries therefore record how to reach that type:
  \[
  \tgt{\begin{aligned}
    Q&::=\kw{bnd}(F)\mid\kw{thru}(F,L)\\[3mm]
    \vec Q&::=[\,]\mid\vec Q,Q.
  \end{aligned}}
  \]
  \gap
  $\tgt{\kw{thru}(F,L)}$ first follows $\tgt{F}$, then applies
  the local template $\tgt{L}$ to the resulting view.
  \gap
  Object entries contain no routes. A route ends at a local template,
  so routes cannot nest as entries.
  \captionline{These separate grammars make the structure used by composition explicit.
  $\tgt{\kw{into}(\vec Q)}$ can start from a type without a source telescope.}
\end{frame}

\begin{frame}{A view contains normalized evidence for each fact}
  \[
  \tgt{\begin{aligned}
    P^\circ&::=\kw{le}(F)\mid\kw{eq}\mid\kw{has}(x,a)\mid\kw{bnd}(F)\\[3mm]
    V&::=[\,]\mid V,P^\circ.
  \end{aligned}}
  \]
  \gap
  For the precise type $\tgt{\obj{z}{[\nm{z}{a}\doteq\rho,\has{a}]}}$,
  the stored literal at $\tgt{x}$ supplies
  \[
    \tgt{V=[\kw{eq},\kw{has}(x,a)]}.
  \]
  \gap
  A cast transforms this view using its normalized recipes.
  Reading a member selects the resulting entry.
  \captionline{An equality entry asserts equal resolved shapes.
  A presence entry identifies a field at an actual location.}
\end{frame}

\begin{frame}{Closed evidence has typed head forms}
  Assume a typed store $\tgt{\vdash\Sigma:\Delta}$.
  \[
    \tgt{\Delta\vdash e:\tau\leq\sigma}
    \quad\Longrightarrow\quad
    \tgt{\exists F.\ N_\Sigma(e)=F\ \land\ \Delta\models F:\tau\leq\sigma}.
  \]
  \captionline{$\tgt{N_\Sigma}$ is head normalization.
  $\tgt{\Delta\models F:\tau\leq\sigma}$ checks the form against resolved endpoint shapes.}
  \gap
  The proof is mutual structural induction on evidence and atom typing.
  It uses typed composition and produces typed views at the same time.
  \gap
  Function forms retain evidence under their parameter binder.
  \captionline{Only evidence closed over the store is normalized here.
  The theorem does not assert termination of programs.}
\end{frame}

\begin{frame}{Evidence recovers shapes in typed stores}
  The target must still prove \textbf{canonical forms} over a typed store:
  \gap
  \[
  \begin{array}{c@{\quad\Longrightarrow\quad}l}
    \tgt{\Delta\vdash r:\Pi(x:\tau)\sigma}
      &\text{a lambda at }\tgt{\rt{r}}\\[4mm]
    \tgt{\Delta\vdash p:x\ni a}
      &\text{field }\tgt{a}\text{ exists at }\tgt{x}
  \end{array}
  \]
  \gap
  Application reaches a lambda. Projection reaches an existing field.
  \captionline{The forms also supply typed argument/result casts.
    This recovers the runtime shapes that direct DOT proofs obtain through
    tight and invertible typing.}
\end{frame}

\begin{frame}{Allocation: keep the precise type at the new location}
  \textbf{Direct proof:} recover inert $\src{P}$ with $\src{\Gamma\vdash P\sub T}$.
  \[
    \src{\Gamma,x:T\vdash t:R}
    \quad\xRightarrow{\text{narrowing}}\quad
    \src{\Gamma,x:P\vdash t:R}.
  \]
  Store the value at $\src{P}$; the continuation syntax stays $\src{t}$.
  \gap
  \textbf{Target:} the literal has type $\tgt{\tau_0}$ and a cast supplies
  $\tgt{\Delta\vdash e:\tau_0\leq\tau}$.
  \[
    \tgt{u':=u[x\mapsto\cast{x}{e}]}.
  \]
  \[
    \tgt{\Delta,x:\tau\vdash u:\rho}
    \quad\xRightarrow{\text{typed substitution}}\quad
    \tgt{\Delta,x:\tau_0\vdash u':\rho}.
  \]
  \captionline{Store the bare literal at $\tgt{\tau_0}$ and continue with $\tgt{u'}$.
  Its transparent binding records witnesses and field labels.}
\end{frame}

\begin{frame}{Application: the direct preservation argument}
  Under $\src{I_D}$, the call has $\src{\Gamma\vdash f:\forall(x:S)T}$ and $\src{\Gamma\vdash y:S}$.
  \gap
  Inversion and store typing recover
  \[
  \src{\begin{gathered}
    \sigma(f)=\lambda(x:S_0).t_0,\qquad\Gamma,x:S_0\vdash t_0:T_0,\\[2mm]
    \Gamma\vdash S\sub S_0,\qquad\Gamma,x:S\vdash T_0\sub T.
  \end{gathered}}
  \]
  \gap
  Narrow the body to $\src{x:S}$; subsume its result to $\src{T}$:
  \[
    \src{\Gamma,x:S\vdash t_0:T}
    \quad\xRightarrow{\text{substitute }\src{y}}\quad
    \src{\Gamma\vdash t_0[y/x]:T[y/x]}.
  \]
  \captionline{The source step is simply $\src{f\,y\longrightarrow t_0[y/x]}$.
  The proof reconstructs why this body meets the advertised type.}
\end{frame}

\begin{frame}{Application: the target executes the recovered evidence}
  Assume $\tgt{\vdash\Sigma:\Delta}$, with $\tgt{\Delta\vdash r:\Pi(x:\tau)\sigma}$ and $\tgt{\Delta\vdash s:\tau}$.
  Canonical forms yield:
  \[
  \tgt{\begin{gathered}
    \Sigma(\rt{r})=\lambda(x:\tau_0).u_0,\qquad
      \Delta,x:\tau_0\vdash u_0:\sigma_0,\\[2mm]
    \Delta\vdash d:\tau\leq\tau_0,\qquad
      \Delta,x:\tau\vdash c:\sigma_0\leq\sigma.
  \end{gathered}}
  \]
  \gap
  The function-coercion case executes
  \[
    \tgt{r\,s\longrightarrow
      \cast{u_0[x\mapsto\cast{s}{d}]}{c[x\mapsto s]}}.
  \]
  \captionline{Cast the argument, substitute into the closure, then cast the result.}
  \[
    \tgt{\rt{\cast{s}{d}}=\rt{s}}
    \quad\Longrightarrow\quad
    \text{body and cast agree at }\tgt{\sigma_0[\rt{s}/x]}.
  \]
\end{frame}

\begin{frame}{Projection: split field existence from result conversion}
  \textbf{Direct ($\src{I_D}$):} from $\src{\Gamma\vdash x:\fld{a}{T}}$,
  recover field body $\src{t_a}$ and type $\src{S}$:
  \[
    \src{\Gamma\vdash t_a[x/z]:S},\qquad\src{\Gamma\vdash S\sub T}.
  \]
  \captionline{Variable and value inversion justify lookup; substitution and subsumption
  type the source reduct.}
  \gap
  \textbf{Target ($\tgt{\vdash\Sigma:\Delta}$):} the translated projection carries
  \[
    \tgt{\Delta\vdash p:x\ni a},\qquad
    \tgt{\Delta\vdash e:\nm{x}{a}\leq\tr{T}}.
  \]
  Presence canonical forms locate the field. Its stored typing gives
  \[
    \tgt{\Delta\vdash u_a[x/z]:\nm{x}{a}}.
  \]
  \captionline{The pending cast $\tgt{e}$ supplies the advertised result type.
  No source field-subtyping inversion is used by the target step.}
\end{frame}

\begin{frame}{Transitivity still needs a composition theorem}
  Over $\tgt{\vdash\Sigma:\Delta}$, well-typed evidence has \textbf{head forms}:
  \[
    \tgt{N_\Sigma(\kw{trans}(e,f))
      =\kw{combine}(N_\Sigma(e),N_\Sigma(f))}.
  \]
  \gap
  Here $\tgt{N_\Sigma}$ denotes head normalization; $\tgt{\models}$ types its forms.
  \[
    \tgt{\frac{\Delta\models F:\tau\leq\sigma\qquad
      \Delta\models G:\sigma\leq\rho}
      {\exists H.\ \kw{combine}(F,G)=H\ \land\ \Delta\models H:\tau\leq\rho}}
  \]
  \gap
  Composition: induction on the sizes of forms.\\[2mm]
  Canonical forms: induction on evidence typing, using composition.
  \captionline{Function forms retain typed argument/result evidence.
  Head normalization does not eliminate transitivity inside that evidence.}
\end{frame}

\begin{frame}{Templates make object composition structural}
  Given $\tgt{\Delta\vdash e:\tau\leq\sigma}$ and
  $\tgt{\Delta\vdash f:\sigma\leq\rho}$, let
  \[
    \tgt{\Phi_\theta=[\has{a},\nm{z}{a}\sqsubseteq\theta]}.
  \]
  \begin{center}
  \begin{tabular}{@{}l@{\qquad}l@{}}
    \textbf{Object conversion}&\textbf{Recipe for fact 1}\\[3mm]
    $\tgt{\obj{z}{\Phi_\tau}\longrightarrow\obj{z}{\Phi_\sigma}}$
      &$\tgt{H_1\longmapsto H_1;e}$\\[3mm]
    $\tgt{\obj{z}{\Phi_\sigma}\longrightarrow\obj{z}{\Phi_\rho}}$
      &$\tgt{H_1\longmapsto H_1;f}$\\[5mm]
    \textbf{Composite}&$\tgt{H_1\longmapsto H_1;(e;f)}$
  \end{tabular}
  \end{center}
  \captionline{$\tgt{H_1}$ names source fact 1; presence at fact 0 is copied.
  Normalize the side proofs first, then compose their forms.}
  \gap
  Crucially, $\tgt{e,f}$ are typed in $\tgt{\Delta}$,
  \textbf{before adding self} $\tgt{z}$.
\end{frame}

\begin{frame}{Why forbid arbitrary evidence under self?}
  An unrestricted object morphism could contain an evidence body
  \[
    \tgt{\Delta,z:\obj{w}{\Phi}\vdash k:\tau\leq\sigma}.
  \]
  To use it at a receiver $\tgt{r}$, one would substitute:
  \[
    \tgt{k':=k[z\mapsto r]}.
  \]
  \begin{center}
  \begin{tabular}{@{}l@{\qquad}l@{}}
    Before substitution & $\tgt{z}$ is hypothetical: no stored literal.\\[3mm]
    After substitution & $\tgt{k'}$ need not be smaller than $\tgt{k}$.
  \end{tabular}
  \end{center}
  \captionline{The closed-store induction gives no hypothesis for this new body,
  which may perform further member eliminations through the receiver.}
  \gap
  Templates reduce instantiation to \textbf{fact lookup and form composition}.
\end{frame}

\begin{frame}{Narrowing-like work becomes cast substitution}
  Suppose
  \[
    \tgt{\Delta\vdash e:\tau'\leq\tau},\qquad
    \tgt{\Delta,x:\tau\vdash f:\sigma\leq\rho}.
  \]
  \gap
  \begin{center}
  \begin{tabular}{@{}l@{\qquad}l@{}}
    \textbf{New proof term}
      &$\tgt{f':=f[x\mapsto\cast{x}{e}]}$\\[5mm]
    \textbf{Typing judgment}
      &$\tgt{\Delta,x:\tau'\vdash f':\sigma\leq\rho}$
  \end{tabular}
  \end{center}
  \gap
  \[
    \tgt{\rt{\cast{x}{e}}=x}
    \quad\Longrightarrow\quad
    \text{names such as }\tgt{\nm{x}{A}}\text{ stay unchanged.}
  \]
  \captionline{Typed substitution is proved before normalization.
  Function-coercion composition uses it to adapt codomain evidence,
  without normalizing under the binder.}
\end{frame}

\begin{frame}{Bad bounds remain possible hypothetically}
  In the hypothetical context
  \[
    \tgt{\Delta_{\kw{bad}}=x:\obj{z}{[
      \top\sqsubseteq\nm{z}{A},\nm{z}{A}\sqsubseteq\bot]}},
  \]
  the two member facts yield
  \[
    \tgt{\Delta_{\kw{bad}}\vdash
      \kw{trans}(\readp{x}{0},\readp{x}{1}):\top\leq\bot}.
  \]
  \gap
  For a \textbf{typed store}, canonical forms instead imply
  \[
    \tgt{\vdash\Sigma:\Delta}
    \quad\Longrightarrow\quad
    \tgt{\neg\exists e.\ (\Delta\vdash e:\top\leq\bot)}.
  \]
  \captionline{Every such $\tgt{e}$ would have a typed head form;
  no form fits these endpoints.
  Consistency follows from store typing and canonical forms.}
\end{frame}

\begin{frame}{Both proofs separate structural lemmas from shape recovery}
  \begin{center}
  \renewcommand{\arraystretch}{0.9}
  \begin{tabular}{@{}c@{\hspace{10mm}}c@{}}
    \textbf{Direct (2017)}&\textbf{FCdot and translation}\\[1mm]
    Narrowing and precise-type facts&Atom substitution and resolution\\
    $\Downarrow$&$\Downarrow$\\
    Tight to invertible&Typed form composition / application\\
    $\Downarrow$&$\Downarrow$\\
    Selection replacement&Precise views and canonical forms\\
    $\Downarrow$&$\Downarrow$\\
    General to tight&Target preservation and progress\\
    $\Downarrow$&$\Downarrow$\\
    Canonical forms&Translation and erasure simulations\\
    $\Downarrow$&$\Downarrow$\\
    Source preservation and progress&Source safety
  \end{tabular}
  \end{center}
  \captionline{Construction order, distinct from the order of use. Both arguments are acyclic.}
\end{frame}

\begin{frame}{The translation preserves typing}
  For well-formed $\src{\Gamma}$:
  \gap
  \[
  \begin{aligned}
    \src{\mathcal{D}\mathrel{::}\Gamma\vdash S\sub T}
      &\Longrightarrow
        \tgt{\tr{\Gamma}\vdash\tr{\mathcal{D}}:\tr{S}\leq\tr{T}}\\[3mm]
    \src{\mathcal{H}\mathrel{::}\Gamma\vdash t:T}
      &\Longrightarrow
        \tgt{\tr{\Gamma}\vdash\tr{\mathcal{H}}:\tr{T}}
  \end{aligned}
  \]
  \captionline{Context well-formedness checks exact literal definitions and distinct labels.
  Ordinary assumptions may contain bad bounds and arbitrary recursive types.}
  \gap
  The translation follows the source derivation.
  \captionline{The source needs no preservation proof at this stage.}
\end{frame}

\begin{frame}{All derivations give the same erased program}
  For every typing derivation $\src{\mathcal{H}\mathrel{::}\Gamma\vdash t:T}$:
  \[
    \underbrace{\er{\tr{\mathcal{H}}}=\er{\src{t}}}_{\text{same erased program}}
  \]
  \gap
  \[
    \er{\tgt{\cast{u}{e}}}=\er{\tgt{u}}.
  \]
  \gap
  \textbf{Coherence:} different derivations of the same term erase identically.
  \captionline{Erasure removes types, type members, and evidence.}
\end{frame}

\begin{frame}{Source safety follows through the shared runtime}
  After a source step, preserve the \emph{existence} of a typed target witness:
  \[
  \begin{array}{c@{\qquad}c@{\qquad}c@{\quad}l}
    \src{s_D}&\src{\longrightarrow}&\src{s_D'}&\text{\src{source}}\\[0mm]
    \downarrow&&\downarrow&\\[0mm]
    q&\longrightarrow&q'&\text{erased runtime}\\[0mm]
    \uparrow&&\uparrow&\\[0mm]
    \tgt{s_F}&\tgt{\longrightarrow^*}&\tgt{s_F'}&\text{\tgt{typed target}}
  \end{array}
  \]
  \captionline{Erase the source step; reflect it to a target run; apply target preservation.}
  \gap
  This establishes $\kw{Sim}(\src{s_D'})$ without retyping $\src{s_D'}$.
  \captionline{For progress, first exhaust the finite sequence of erased cast steps;
  target progress then yields a source step or a final state.}
\end{frame}

\begin{frame}{DOT-MNF machine safety is a corollary}
  Every state reachable from a closed well-typed program is final or can step:
  \[
  \begin{gathered}
    \src{\varnothing\vdash t:T,\quad\langle\varnothing,\varnothing,t\rangle\longrightarrow^*s}\\[3mm]
    \Longrightarrow\quad\src{\kw{Final}(s)\ \lor\ \exists s'.\,s\longrightarrow s'}
  \end{gathered}
  \]
  \captionline{Final: a value or variable with no pending continuation.}
  \gap
  Shape reasoning uses normalized target evidence.\\[2mm]
  The invariant is a typed target state with the same erasure.
  \captionline{No separate source preservation proof.
  Target substitution, composition, and canonical forms carry the shape reasoning.}
\end{frame}

\begin{frame}{Retained-let reduction keeps bindings in the term}
  The paper's semantics retains allocated values:
  \[
    \src{\kw{let}\ x=v\ \kw{in}\ t}.
  \]
  Write $\src{\sigma\vdash t\rightsquigarrow u}$ for reduction under
  the surrounding value bindings $\src{\sigma}$.
  \[
  \src{\begin{aligned}
    x\,y&\rightsquigarrow t[y/z]
      &&\text{if }\sigma(x)=\lambda(z:S).t\\[3mm]
    x.a&\rightsquigarrow t[x/z]
      &&\text{if }\sigma(x)=\nu(z:T)d,\ a=t\in d\\[3mm]
    \kw{let}\ x=y\ \kw{in}\ t&\rightsquigarrow t[y/x].
  \end{aligned}}
  \]
  \captionline{Projection uses relational membership in the definitions.
  Distinct labels make this agree with deterministic machine lookup.}
\end{frame}

\begin{frame}{Retained reduction allows overlapping orders}
  Both evaluation contexts are available:
  \[
    \src{\frac{\sigma\vdash t\rightsquigarrow t'}
      {\sigma\vdash\kw{let}\ x=t\ \kw{in}\ u
        \rightsquigarrow\kw{let}\ x=t'\ \kw{in}\ u}}
  \]
  \[
    \src{\frac{\sigma[x\mapsto v]\vdash t\rightsquigarrow t'}
      {\sigma\vdash\kw{let}\ x=v\ \kw{in}\ t
        \rightsquigarrow\kw{let}\ x=v\ \kw{in}\ t'}}.
  \]
  Reassociation can also apply:
  \[
  \src{\begin{aligned}
    \kw{let}\ x=(\kw{let}\ y=t\ \kw{in}\ u)\ \kw{in}\ w
      \quad\rightsquigarrow\quad{}\\[-1mm]
    \kw{let}\ y=t\ \kw{in}\ \kw{let}\ x=u\ \kw{in}\ w.
  \end{aligned}}
  \]
  \captionline{There is no priority between reassociation and reduction inside the right-hand side.
  Bound variables are renamed to avoid capture.}
\end{frame}

\begin{frame}{Readback reconstructs retained lets}
  Readback replaces the machine's store and continuation by enclosing lets:
  \[
    \src{\kw{rb}(\langle\sigma,K,t\rangle)
      :=\kw{close}_{\sigma}(\kw{plug}_{K}(t))}.
  \]
  For example,
  \[
    \src{\langle[x\mapsto v],[y.u],t\rangle}
      \quad\longmapsto\quad
    \src{\kw{let}\ x=v\ \kw{in}\ \kw{let}\ y=t\ \kw{in}\ u}.
  \]
  \gap
  Every machine step gives a finite retained-let reduction:
  \[
    \src{s\longrightarrow s'}\quad\Longrightarrow\quad
    \src{\kw{rb}(s)\rightsquigarrow^*\kw{rb}(s')}.
  \]
  \captionline{Pushing a frame leaves readback unchanged.
  Allocation may reassociate a value binding across several frames.}
\end{frame}

\begin{frame}{Answers retain all enclosing value bindings}
  The retained-let answers are
  \[
    \src{q::=x\mid v\mid\kw{let}\ x=v\ \kw{in}\ q}.
  \]
  A stuck term is neither an answer nor able to reduce.
  \gap
  A final machine state reads back to an answer, including its store:
  \[
    \src{\langle[x\mapsto v],\varnothing,x\rangle}
      \quad\longmapsto\quad
    \src{\kw{let}\ x=v\ \kw{in}\ x}.
  \]
  \captionline{Observations are exact annotated terms with retained bindings.
  No quotient by garbage collection or binding reordering is used.}
\end{frame}

\begin{frame}{The converse preserves finite behavior backwards}
  Assume distinct labels throughout $\src{\sigma}$ and $\src{t}$,
  and $\src{\sigma\vdash t\rightsquigarrow u}$.
  \[
    \src{\frac{\langle\sigma,K,u\rangle\Downarrow o}
      {\langle\sigma,K,t\rangle\Downarrow o}}
    \qquad\text{for every continuation }\src{K}.
  \]
  \captionline{$\src{s\Downarrow o}$ means a finite machine run observing $\src{o}$:
  an exact retained answer, or stuckness.}
  \gap
  \begin{center}
  \begin{tabular}{@{}l@{\qquad}l@{}}
    \textbf{Retained step}&\textbf{Machine argument}\\[3mm]
    Reassociation&Fuse pending let frames\\[3mm]
    Under a value binding&Allocate that value
  \end{tabular}
  \end{center}
  \captionline{Typing supplies label distinctness; reduction preserves it.
  No source type preservation or inert-context theorem is needed.}
\end{frame}

\begin{frame}{Exact finite-answer and stuckness correspondence}
  For a closed term $\src{t}$ with distinct object labels throughout,
  and a retained answer $\src{q}$:
  \[
  \src{t\rightsquigarrow^*q}
  \quad\Longleftrightarrow\quad
  \src{\exists s.\ \kw{init}(t)\longrightarrow^*s
    \land\kw{Final}(s)\land\kw{rb}(s)=q}.
  \]
  \gap
  Stuckness also agrees:
  \[
  \src{\exists u.\ t\rightsquigarrow^*u\land\kw{Stuck}(u)}
  \quad\Longleftrightarrow\quad
  \src{\exists s.\ \kw{init}(t)\longrightarrow^*s\land\kw{Stuck}(s)}.
  \]
  \captionline{These theorems cover every finite retained reduction order.
  The same exact answer appears on both sides of the first equivalence.}
\end{frame}

\begin{frame}{Safety for every retained-let reduction order}
  For every closed annotated typing derivation and every retained run:
  \[
    \src{\varnothing\vdash_{\!W}t:T},\qquad
    \src{t\rightsquigarrow^*u}
    \quad\Longrightarrow\quad
    \src{\kw{Answer}(u)\ \lor\ \exists v.\ u\rightsquigarrow v}.
  \]
  \gap
  A stuck retained reduct would give a stuck annotated machine execution.
  Annotation erasure and the FCdot translation rule out that execution.
  \gap
  The proof therefore covers all choices of retained reduction order.
  \captionline{The source contributes label distinctness and operational correspondence.
  The target contributes preservation, canonical forms, and progress.}
\end{frame}

\begin{frame}{Mechanized scope}
  The public source distinguishes type labels $\src{A,B}$ from field labels
  $\src{a,b}$ throughout syntax and typing derivations.
  \gap
  It admits arbitrary recursive bodies, unrestricted intersections,
  and same-block aliases.
  \gap
  Paths are variables. Scoping is intrinsic; the slides use named notation
  for the binders.
  \captionline{A conversion from named syntax modulo alpha-equivalence is not formalized.
  Stable field paths and recursive subtyping require separate extensions.}
\end{frame}

\appendix
\begin{frame}{Reference: key DOT-MNF rules}
  \[
    \src{\frac{\Gamma\vdash t:S\quad\Gamma\vdash S\sub T}{\Gamma\vdash t:T}}
    \qquad
    \src{\frac{\Gamma\vdash x:\fld{a}{T}}{\Gamma\vdash x.a:T}}
  \]
  \[
    \src{\frac{\Gamma\vdash x:\mu(z.T)}
      {\Gamma\vdash x:T[x/z]}}
    \qquad
    \src{\Gamma\vdash(A=T):\mem{A}{T}{T}}
  \]
  \captionline{Recursive introduction and elimination accept every body $\src{T}$.
  A type definition supplies equal lower and upper bounds.}
\end{frame}

\begin{frame}{Reference: inclusion, equality, and presence evidence}
  \[
  \tgt{\begin{aligned}
    e,f&::=\kw{refl}(\tau)\mid\kw{trans}(e,f)\mid\kw{top}(\tau)\mid\kw{bot}(\tau)\\[1mm]
      &\quad\mid\kw{eqToLe}(\varphi)\mid\kw{pi}(e,x.f)\mid\kw{obj}(\Phi,m)\\[1mm]
      &\quad\mid\kw{pair}(\Phi,\Psi,e,f)\mid\kw{bound}(\Phi,i)\\[1mm]
      &\quad\mid\kw{intoBnd}(e)\mid\kw{member}(r,e,i)\\[3mm]
    \varphi&::=\kw{refl}(\tau)\mid\kw{symm}(\varphi)
      \mid\kw{trans}(\varphi,\varphi)\\[1mm]
      &\quad\mid\kw{def}(x,\ell)\mid\kw{member}(r,e,i)\\[3mm]
    p&::=\kw{field}(a)\mid\kw{member}(r,e,i)
  \end{aligned}}
  \]
  \captionline{The parameter $\tgt{x}$ scopes over the codomain evidence in $\tgt{\kw{pi}(e,x.f)}$.
  A telescope annotation has its own self binder.}
\end{frame}

\begin{frame}{Reference: object morphisms and their judgments}
  \[
  \tgt{\begin{aligned}
    H&::=\kw{le}(j)\mid\kw{eq}(j)\mid\kw{eqSym}(j)\\[2mm]
    q&::=\kw{none}\mid\kw{some}(e)\\[2mm]
    m&::=\kw{nil}\mid\kw{le}(m,q,H,q)\mid\kw{eq}(m,j,b)\\[2mm]
      &\quad\mid\kw{has}(m,j)\mid\kw{bnd}(m,e)
  \end{aligned}}
  \]
  \captionline{$\tgt{H}$ names a fact; $\tgt{q}$ is an optional side;
  $\tgt{b}$ reverses equality. Present sides are checked before adding self.}
  \gap
  \[
  \begin{array}{cc}
  \tgt{\Delta\vdash e:\tau\leq\sigma}&
  \tgt{\Delta\vdash\varphi:\tau\equiv\sigma}\\[2mm]
  \tgt{\Delta\vdash p:x\ni a}&
  \tgt{\Delta\vdash m:\Phi\Rightarrow\Psi}
  \end{array}
  \]
\end{frame}

\begin{frame}{Reference: the two telescope translations}
  Let $\src{D}$ range over top, declarations, their intersections,
  and recursively declaration-shaped types.
  \[
  \begin{aligned}
    \tel{\top}&=\tgt{[\,]}\qquad
    \tel{S\wedge T}=\tgt{\tel{S}\mathbin{+\!+}\tel{T}}\\[2mm]
    \tel{\mem{A}{S}{T}}&=\tgt{[\tr{S}\sqsubseteq\nm{z}{A},\nm{z}{A}\sqsubseteq\tr{T}]}\\[2mm]
    \tel{\fld{a}{T}}&=\tgt{[\has{a},\nm{z}{a}\sqsubseteq\tr{T}]}\\[2mm]
    \tel{\mu(z.D)}&=\tels{z}{D}\\[2mm]
    \tel{B}&=\tgt{[\bnd{\tr{B}}]}\qquad\text{for every other shape }\src{B}.
  \end{aligned}
  \]
  \captionline{In $\tel{T}$, translated endpoints are weakened under fresh self $\tgt{z}$.
  $\tels{z}{T}$ uses the body's existing self instead, so its bounds can mention $\src{z}$.
  These shape cases organize translation and do not restrict source typing.}
\end{frame}

\begin{frame}{Reference: recursive adapters and precise types}
  For a source body $\src{T}$, root $\tgt{x}$, and translated atom $\tgt{r}$:
  \[
    \tgt{\kw{recI}(r)=\kw{fold}(z.\tels{z}{T},
      \kw{unfold}(\kw{into}_{\src{T[x/z]}}(r)))}.
  \]
  \[
    \tgt{r':=\kw{fold}(w.\tel{T[x/z]},\kw{unfold}(r))}.
  \]
  Recursive elimination returns $\tgt{r'}$, extracting its single bound
  when the opened body is not an object shape.
  \captionline{$\tgt{\kw{into}_{\src{U}}}$ is the identity for object shapes.
  Otherwise it casts into $\tgt{\obj{w}{[\bnd{\tr{U}}]}}$.
  Every wrapper preserves the root.}
  \[
    \tgt{W=[a=\tau],\quad F=[a=u]}
    \quad\Longrightarrow\quad
    \tgt{\kw{precise}(W,F)=\obj{z}{[\nm{z}{a}\doteq\tau,\has{a}]}}.
  \]
  \captionline{In general: an equality per witness, presence per field.
  A typed store keeps bare literals at these types; casts remain at their uses.}
\end{frame}

\begin{frame}{Reference: pushback is a different direct strategy}
  \href{https://namin.seas.harvard.edu/files/soundness_oopsla16.pdf}{Rompf--Amin (2016)}
  study a larger DOT with recursive subtyping.
  \[
  \begin{array}{c}
    \text{primitive transitivity supports structural narrowing}\\[1mm]
    \Downarrow\\[1mm]
    \text{pushback in concrete runtime contexts}\\[1mm]
    \Downarrow\\[1mm]
    \text{subtyping inversion}
  \end{array}
  \]
  \captionline{Pushback is required only when the abstract context is empty.}
  \gap
  Our source has variable Rec-I/Rec-E and no recursive subtyping rule.
  \captionline{The main comparison uses Rapoport et al. (2017),
  whose proof instead passes through tight and invertible typing.}
\end{frame}

\end{document}
